\section{Beliefs, rationality and strategic choice}
\label{sec:annex-p17-epistemic-game-theory}

\begin{annexitem}{annex:p17.egt.hierarchies}{thm:p17-egt-hierarchies}{Finite pushforwards and consistent truncation}
\begin{proof}
For any set \(X\), let \(\Delta_f(X)\) denote the probabilities
with finite support, defined by their nonnegative atom weights
summing to one. A map \(f:X\to Y\) pushes such a probability
forward by summing the weights of atoms with the same image.
The result has finite support and total mass one. Successive
pushforwards compose, since both procedures sum precisely the
weights of the same original atoms.

Define \(\mathcal H_i^0=\{*\}\) and
\[
 \mathcal H_i^{r+1}
   =\Delta_f\!\left(\prod_{j\ne i}(C_j\times\mathcal H_j^r)\right).
\]
The finite source law \(b_i(t_i)\) and induction make
\(h_i^{r+1}(t_i)\) the pushforward into this space under
\((c_j,t_j)_j\mapsto(c_j,h_j^r(t_j))_j\).
This defines each level without choosing representatives for
types that happen to have identical lower-level beliefs.
Their probabilities are simply added.

Let \(\pi_i^0:\mathcal H_i^1\to\mathcal H_i^0\) be the constant
map. For \(r\geq1\), let \(\pi_i^r\) push a distribution in
\(\mathcal H_i^{r+1}\) forward under
\[
 (c_j,\eta_j)_j\longmapsto(c_j,\pi_j^{r-1}(\eta_j))_j.
\]
For \(r=0\), \(\pi_i^0h_i^1=h_i^0\).
If \(\pi_j^{r-1}h_j^r=h_j^{r-1}\) for every opponent,
composition of the finite pushforwards gives
\(\pi_i^rh_i^{r+1}=h_i^r\).
Induction proves every consistency identity.
For one player the opponent product is a singleton at every
stage, so these same constructions remain defined.
The infinite sequence is specified by one value for each natural
depth; no universal representation theorem has been assumed.
\end{proof}
\end{annexitem}

\begin{annexitem}{annex:p17.egt.belief.algorithm}{thm:p17-egt-belief-algorithm}{Cumulative levels and reachability of a failure}
\begin{proof}
For any family of sets of opponent types, probability one of
membership means that every positive-probability tuple has all
its coordinates in the respective sets. For finitely many such
requirements this is equivalent to membership in their intersection.
Using the definition of \(H_i^{r+1}\) therefore gives
\[
 \bigcap_{\ell=1}^{r+1}H_i^\ell
   =H_i^1\cap\{t_i:\text{all successors are in }
                         \bigcap_{\ell=1}^{r}H_j^\ell\}.
\]
This proves the displayed recursion, including \(r=0\).
It also shows \(Q_i^{r+1}\subseteq Q_i^r\) by induction.

A vertex is in \(Q_i^1\) exactly when it is not itself bad.
Inductively, the recursion says that it is in \(Q_i^r\)
exactly when no bad vertex can be reached by a path of length
less than \(r\). An edge records positive marginal probability;
in a finite distribution such a marginal is positive exactly
when some full tuple containing that type has positive mass.
Thus the graph and tuple conditions agree.

If a bad vertex is reachable, erase loops from a witnessing path.
With \(N=\sum_i|T_i|\), the resulting path has length at most
\(N-1\). Consequently \(Q^N=Q^\infty\). Alternatively, any
nonstationary removal step discards at least one vertex, and a
stationary step stays stationary under the same recursion.
This proves the finite bound and the reachability criterion.

For failure of monotonicity of the individual levels, give both
players choices \(0,1\) and utility equal to their own choice.
Choice \(1\) is uniquely optimal at every type.
Take \(T_1=\{x,z\}\), \(T_2=\{y,w\}\), with
\[
 b_1(x)=(0,y),\quad b_1(z)=(1,w),\quad
 b_2(y)=(1,z),\quad b_2(w)=(1,z).
\]
Type \(x\) fails onefold belief because it assigns certainty to
the opponent's nonoptimal choice \(0\).
Type \(y\) expresses onefold belief because it assigns certainty
to the rational pair \((1,z)\).
Hence \(x\) expresses twofold belief. The types \(z,w\) form a
closed cycle of rational pairs, and \(y\) points into that cycle,
so they express all levels. In fact \(x\) expresses every
individual level from two onward but never common belief.
This verifies why the cumulative intersection cannot be omitted.
\end{proof}
\end{annexitem}

\begin{annexitem}{annex:p17.egt.best.response.alternative}{thm:p17-egt-best-response-alternative}{The complete finite alternative}
\begin{proof}
Use the real matrix \(D\) from the statement as a zero-sum
matrix game. Its rows are \(R_i\); its columns are the finite
nonempty set \(R_{-i}\), including the singleton empty tuple
when there are no opponents.
By \cref{thm:p17-finite-minimax}, both extrema are attained and
\[
 v=\max_{x\in\Delta(R_i)}\min_{c_{-i}}x^{\mathsf T}D e_{c_{-i}}
   =\min_{p\in\Delta(R_{-i})}\max_{d_i}(Dp)_{d_i}.
\]
The row for \(a_i\) is identically zero. Therefore \(v\geq0\).
If \(v>0\), a maximizing \(x\) has every coordinate of
\(x^{\mathsf T}D\) at least \(v>0\), providing a strict
domination witness. If \(v=0\), a minimizing \(p\) has
\((Dp)_{d_i}\leq0\) at every row, making \(a_i\) a best response.
These exhaust the possibilities for \(v\).

They cannot both hold: multiplying \(Dp\leq0\) by \(x\geq0\)
would give \(x^{\mathsf T}Dp\leq0\), whereas a strictly positive
row vector \(x^{\mathsf T}D\) averaged by a probability \(p\)
has positive value. Conversely, each of the displayed
inequalities has exactly the payoff meaning specified in
the statement. This proves both directions and exclusivity.

For the exact rational computation, put \(M=\max|D_{d,c}|\).
The margin program with \(x\) in its simplex and \(0\leq z\leq M\)
has a nonempty compact polytope of feasible points: the zero
row's unit mass with \(z=0\) is feasible. Its maximum equals
the displayed \(v\), since \(0\leq v\leq M\).
The belief polytope is nonempty when \(v=0\), by the argument above,
and is compact as a closed subset of a probability simplex.

Here is the needed vertex argument, including lower-dimensional
polytopes. Express a nonempty bounded polytope by finitely many
linear inequalities and equalities. A linear objective attains
its maximum by compactness. Among its maximizing points choose
one with the largest rank of active constraint normals, counting
the equality normals throughout. If that rank is smaller than
the ambient variable dimension, choose a nonzero vector annihilated
by all active normals. Small displacements in either direction
preserve the active constraints and, by the positive slack of each
remaining inequality, remain feasible. Maximality of the objective
then forces its change in this direction to be zero.
Move along the positive ray until a new inequality becomes tight.
Such an inequality must exist because the polytope is bounded;
the first endpoint is the minimum of finitely many positive
slack-to-slope ratios. The new normal has a nonzero value on
the direction and is therefore independent of the old normals.
The endpoint is still a maximizer, contradicting maximal rank.
Thus a full-rank active system determines a maximizing point.
Applying the same reasoning with zero objective proves the
existence of such a point in every nonempty bounded polytope.

There are finitely many full-rank subsets of the given constraint
normals. Solve each associated equality system, reject infeasible
solutions, and compare objective values among those retained.
Every solution is rational for rational data by finite Gaussian
elimination, and the preceding paragraph guarantees inclusion of
an optimum. For the belief polytope choose any feasible vertex.
This proves the finite exact procedure, also when \(M=0\), when
some probabilities vanish, or when the feasible set is a singleton.
\end{proof}
\end{annexitem}

\begin{annexitem}{annex:p17.egt.reduction.integrity}{thm:p17-egt-reduction-integrity}{Retaining maxima and preserving all original deviations}
\begin{proof}
Assume \(R^k\) is nonempty. Choose any probability on
\(R_{-i}^k\), and maximize its expected payoff on \(R_i^k\).
A maximizing pure action cannot be strictly dominated there:
averaging a strict domination inequality would contradict
its maximality. Thus \(R_i^{k+1}\ne\varnothing\).
Starting with the given nonempty sets proves nonemptiness
for every round and every player.

We prove the stronger assertion that, for any probability \(p\)
on \(R_{-i}^k\), every action outside \(R_i^k\) has payoff
strictly below the maximum over \(R_i^k\). It is vacuous at zero.
Suppose it holds at \(k\), and take \(p\) on
\(R_{-i}^{k+1}\subseteq R_{-i}^k\).
An action removed from \(R_i^k\) in this round is strictly
dominated on \(R_{-i}^k\) by a mixture on \(R_i^k\).
Its \(p\)-payoff is therefore strictly below the maximum
over \(R_i^k\). No action attaining that maximum is removed,
so this maximum equals the maximum over \(R_i^{k+1}\).
The induction hypothesis already puts every earlier-removed
action strictly below the former, and hence the latter, maximum.
This proves the assertion at \(k+1\).

In particular, all original maximizers are retained, and maximizing
over the retained actions has the same value as maximizing over
the original action set for the specified beliefs. If \(a_i\)
is retained at round \(k+1\), its belief witness in the restricted
game \(R^k\), obtained from the preceding lemma, therefore makes
it optimal over \(C_i\), too.

The sets decrease and never become empty. Each strict round
deletes at least one action, while at least one per player remains.
There can be at most \(\sum_i(|C_i|-1)\) such rounds.
If a round deletes nothing, applying the same rule to the unchanged
game deletes nothing subsequently. This proves stabilization.
\end{proof}
\end{annexitem}

\begin{annexitem}{annex:p17.egt.finite.depth}{thm:p17-egt-finite-depth}{Both directions and an explicit finite type construction}
\begin{proof}
First prove, by induction on \(k\), that a choice optimal for a
type in \(Q_i^k\) survives \(k+1\) rounds in the original game.
For \(k=0\), its first-order belief makes it undominated in
the full game by the belief-or-domination alternative.
For \(k\geq1\), a type in \(Q_i^k\) assigns positive probability
only to rational opponent pairs \((c_j,t_j)\) with
\(t_j\in Q_j^{k-1}\). The induction hypothesis puts each such
\(c_j\) in \(R_j^k\). Its first-order belief is consequently
supported on \(R_{-i}^k\). The choice itself belongs to \(R_i^k\),
by the preceding-depth assertion, and is optimal against
all original deviations, hence against those in \(R_i^k\).
The alternative applied to \(R^k\) retains it in \(R_i^{k+1}\).

For the converse fix a finite depth \(k\). For every player \(i\),
level \(0\leq m\leq k\), and action \(a\in R_i^{m+1}\),
choose a witness
\[
 p_i^{m,a}\in\Delta(R_{-i}^m)
\]
making \(a\) optimal in \(R_i^m\). There are finitely many
such choices. The reduction-integrity theorem makes every
one of these witnesses optimal over the original \(C_i\).
Introduce a type \(t_i^{m,a}\) for each witness, and a further
anchor type \(\tau_i\). Choose one arbitrary \(c_i^0\in C_i\).
Define the anchor beliefs by
\[
             b_i(\tau_i)=\delta_{(c_j^0,\tau_j)_{j\ne i}}.
\]
At level zero push \(p_i^{0,a}\) forward by
\[
             c_{-i}\longmapsto(c_j,\tau_j)_{j\ne i}.
\]
At level \(m\geq1\), instead push \(p_i^{m,a}\) forward by
\[
             c_{-i}\longmapsto(c_j,t_j^{m-1,c_j})_{j\ne i}.
\]
All named types exist because \(c_j\in R_j^m\) on the
support of that probability. These definitions specify every
belief of a finite nonempty epistemic model.

The first-order marginal at \(t_i^{m,a}\) is exactly its
witness \(p_i^{m,a}\), so \(a\) is optimal for it.
At level zero no belief-in-rationality condition is required.
Inductively, a level-\(m\) type believes in the rationality
of every supported opponent pair, and each of those opponent
types expresses all levels through \(m-1\).
The cumulative recursion therefore puts this type in \(Q_i^m\).
In particular \(t_i^{k,a_i}\) supplies the requested witness.
\end{proof}
\end{annexitem}

\begin{annexitem}{annex:p17.egt.fundamental}{thm:p17-egt-fundamental}{The full finite static equivalence}
\begin{proof}
If \(a_i\) is optimal at a common-belief type, then that type
lies in \(Q_i^k\) at every finite depth. The preceding theorem
places \(a_i\) in \(R_i^{k+1}\) at every depth, hence in
\(R_i^\infty\).

Conversely let \(R^\infty=R^K\) be the nonempty stationary
restriction. Every \(a\in R_i^\infty\) is undominated in that
restriction and therefore has a belief \(p_i^a\) supported on
\(R_{-i}^\infty\) making it optimal over \(R_i^\infty\).
Reduction integrity makes it optimal over \(C_i\).
Set \(T_i=\{t_i^a:a\in R_i^\infty\}\) and push \(p_i^a\)
forward by \(c_{-i}\mapsto(c_j,t_j^{c_j})_{j\ne i}\).
Every conjectured choice is optimal for its conjectured type,
so every type has onefold belief. Every successor is another
such type, so induction gives every higher level at every type.
The model is finite, all types have common belief, and its
type \(t_i^{a_i}\) rationalizes the selected action.
This establishes the claimed existential equivalence on the
entire finite game, with arbitrary correlated conjectures.
\end{proof}
\end{annexitem}

\begin{annexitem}{annex:p17.egt.first.game}{thm:p17-egt-first-game}{All computations in the four-by-four source example}
\begin{proof}
In action order \(a,b,c,d\), the expected-payoff vectors for
the three player-1 types are
\[
 v_1(\cdot,t_1^a)=(4,3,0,1),\quad
 v_1(\cdot,t_1^b)=(0,3,2,1),\quad
 v_1(\cdot,t_1^c)=(8/5,9/5,2,1).
\]
In action order \(e,f,g,h\), the player-2 vectors are
\[
 v_2(\cdot,t_2^e)=(2,1,8/5,9/5),\quad
 v_2(\cdot,t_2^g)=(0,1,4,3),\quad
 v_2(\cdot,t_2^h)=(2,1,0,3).
\]
Each vector has the asserted unique maximum.
Type \(t_1^c\) assigns positive mass to \((f,t_2^g)\),
where \(f\) is not optimal; \(t_2^e\) assigns positive
mass to \((d,t_1^a)\), where \(d\) is not optimal.
Every other conjectured pair in the six assignments is rational.
These are exactly the two bad types.

The type \(t_1^b\) points to the bad \(t_2^e\), and
\(t_2^h\) points to the bad \(t_1^c\), so these two types
fail the two-level cumulative condition.
The types \(t_1^a,t_2^g\) form a closed rational cycle and
therefore have common belief. No other type can have it
because it is bad or reaches a bad vertex.
For example, \(t_1^b\)'s first-order belief assigns certainty
to \(e\); its second-order belief assigns certainty to the
opponent's first-order probability \(3c/5+2d/5\).
This computes the source's hierarchy example with the
opponent type carrying player index \(2\).

For elimination, the mixture \((a+b)/2\) has player-1 payoffs
\((3/2,2,7/2,7/2)\), strictly above \(d\)'s
\((1,1,1,0)\). The player-2 mixture \((g+h)/2\) has payoffs
\((7/2,7/2,3/2,2)\) across rows, strictly above \(f\)'s
\((1,0,1,1)\).
None of \(a,b,c,e,g,h\) is dominated in the original game:
the six type marginals above are their explicit best-response
witnesses. This proves exactly the first round.

In the restriction \(\{a,b,c\}\times\{e,g,h\}\),
\(b\)'s vector \((3,3,3)\) strictly exceeds \(c\)'s
\((2,0,2)\), while \(h\)'s \((3,3,3)\) strictly exceeds
\(e\)'s \((0,2,2)\).
Choices \(a,b\) remain best responses to \(g,e\),
respectively; choices \(g,h\) remain best responses to
\(a,c\), respectively. Those witnesses are in the current
restriction before the simultaneous deletions.
Thus the second restriction is exactly
\(\{a,b\}\times\{g,h\}\).

There \(a\) gives \(4\) against both columns, exceeding
\(b\)'s \(3\), and \(g\) gives \(4\) against both rows,
exceeding \(h\)'s \(3\). The third restriction consists
of \(a,g\), and a singleton action set has no strict
domination of its sole action. This completes every round.
\end{proof}
\end{annexitem}

\begin{annexitem}{annex:p17.egt.nash.beliefs}{thm:p17-egt-nash-beliefs}{The simple equilibrium hierarchy}
\begin{proof}
The product of the finitely many opponent distributions
is a probability distribution and defines each \(b_i\).
Its first-order marginal is the product law
\(\bigotimes_{j\ne i}\sigma_j\).
By the equilibrium support conditions, every \(c_j\) with
\(\sigma_j(c_j)>0\) is a best response to
\(\sigma_{-j}\), precisely the marginal at the sole type
\(t_j\). Hence every positive-probability opponent tuple
contains only rational choice-type pairs.
All types have onefold belief, and all successors remain
in this same set of types. Induction gives common belief.
The same equilibrium support condition proves optimality
of each supported own choice at its type.
The result specifies probability laws on conjectured choices;
it adds no claim about the physical mechanism of actual choice.
\end{proof}
\end{annexitem}

\begin{annexitem}{annex:p17.egt.nash.gap}{thm:p17-egt-nash-gap}{The six-type cycle and uniqueness of the mixed equilibrium}
\begin{proof}
Against \(e,d,f\), player 1's uniquely optimal choices are
\(a,b,c\), respectively, with payoffs \(3,3,2\).
Against \(a,b,c\), player 2's uniquely optimal choices are
\(d,e,f\), respectively, with payoffs \(3,3,2\).
Every arrow in the specified model therefore reaches a
rational pair. The four types
\(t_1^a,t_2^e,t_1^b,t_2^d\) form a cycle; the other two
form a second cycle. The graph criterion gives common belief
at every type and rational possibility of every action.

For a completely general mixed profile, write player 1's
probabilities as \((x,y,z)\) on \(a,b,c\) and player 2's
as \((p,q,r)\) on \(d,e,f\).
Their pure-action payoff vectors are
\[
                  (3q,3p,2),\qquad(3x,3y,2).
\]
If \(x>0\) in equilibrium, optimality of \(a\) requires
\(q\geq2/3\). Then \(q>0\) requires \(y\geq2/3\);
then \(y>0\) requires \(p\geq2/3\), contradicting
\(p+q\leq1\). If \(y>0\), the same cycle in the other
direction gives \(p\geq2/3\), \(x\geq2/3\), and
\(q\geq2/3\), again a contradiction.
Thus \(x=y=0\), \(z=1\), after which the second player's
unique optimum is \(f\), giving \(p=q=0,r=1\).
Against \(f\), the first player's unique optimum is \(c\).
This verifies existence and uniqueness of the asserted equilibrium.

Finally the first three arrows from \(t_1^a\) are
\((e,t_2^e)\), then \((b,t_1^b)\), then \((d,t_2^d)\).
They give exactly the stated discrepancy between the
first-order belief \(e\) and the higher conjecture about
that belief, namely \(d\).
\end{proof}
\end{annexitem}

\begin{annexitem}{annex:p17.egt.mutual.conjectures}{thm:p17-egt-mutual-conjectures}{Actual conjectures and mutual correctness}
\begin{proof}
Put \(\sigma_1=q_2(t_2)\) and \(\sigma_2=q_1(t_1)\).
If \(\sigma_2(c_2)>0\), finite marginalization supplies
some \(s_2\) with \(b_1(t_1)(c_2,s_2)>0\).
Belief in rationality makes \(c_2\) optimal at \(s_2\).
The correctness hypothesis gives
\(q_2(s_2)=q_2(t_2)=\sigma_1\).
Thus every action in the support of \(\sigma_2\) is a
best response to \(\sigma_1\).
Interchanging the two players proves that every supported
action of \(\sigma_1\) is a best response to \(\sigma_2\).
These are exactly the finite Nash support conditions.

Both probability-one requirements are relative to the
specified actual types. No common prior or common belief
closure is needed for this implication. Nor has it been
assumed that a player's realized action receives positive
probability in the other player's conjecture: the conclusion
identifies the conjecture profile, not an additional support
condition on unmentioned realized actions.
\end{proof}
\end{annexitem}

\begin{annexitem}{annex:p17.egt.many.player.boundaries}{thm:p17-egt-many-player-boundaries}{Two full three-player countermodels}
\begin{proof}
For correlation take \(C_1=\{A,B\}\), \(C_2=C_3=\{0,1\}\).
Give player 1 payoff \(1\) for \(A\) when the other choices
agree and \(0\) otherwise, and constant payoff \(3/4\) for \(B\).
Give the other players payoff zero everywhere.
There is one type for each player. Type 1 assigns probability
one-half to each opponent profile \((0,0)\), \((1,1)\).
Type 2 assigns certainty to player 1 choosing \(A\) and
a fair distribution to player 3's choice; type 3 assigns
certainty to \(A\) and a fair distribution to player 2.
Attach the sole type to every conjectured choice.

At type 1, the payoffs of \(A,B\) are \(1,3/4\), so \(A\)
is uniquely optimal. Both choices of each other player are
optimal. Thus all conjectured pairs are rational and every
type has common belief. All players who conjecture about
the same player agree on that player's marginal, but type 1's
opponent law is correlated. Replacing it by the product of
its two fair marginals makes \(A\)'s payoff \(1/2\), so \(B\)
becomes uniquely optimal. In particular the common marginals
\((\delta_A,\mathrm{fair},\mathrm{fair})\) do not themselves
form a Nash equilibrium.

For disagreement give all three players choices \(0,1\)
and utility zero at every profile. Again give each one type.
Type 1 assigns certainty to opponents \((c_2,c_3)=(0,0)\);
type 2 to \((c_1,c_3)=(0,1)\); and type 3 to
\((c_1,c_2)=(0,0)\). All these opponent laws are products.
All choices are rational, so all types have common belief.
Nevertheless players 1 and 2 assign different point masses
to player 3's action. No single conjecture \(\sigma_3\)
equals both. This proves that independence does not also
supply agreement.
For any larger finite player set, append players with one action,
zero payoff and one type, and extend every old belief by their
certain choice-type coordinates. The same two counterexamples
and all common-belief conclusions are then preserved.
\end{proof}
\end{annexitem}

\begin{annexitem}{annex:p17.egt.reny.reasoning}{thm:p17-egt-reny-reasoning}{Every contingent payoff and both conditional reasoning paths}
\begin{proof}
The complete strategic-form table implied by the five terminal
histories is
\[
\begin{array}{c|rrrr}
 &cg&ch&dg&dh\\\hline
 ae&(3,3)&(3,3)&(3,3)&(3,3)\\
 af&(3,3)&(3,3)&(3,3)&(3,3)\\
 be&(2,2)&(2,2)&(1,1)&(1,1)\\
 bf&(2,2)&(2,2)&(0,4)&(4,0).
\end{array}
\]
Backward at the last node, player 2 prefers \(g\) since \(4>0\).
Given that continuation, player 1 prefers \(e\) since \(1>0\).
Player 2 at the earlier node then prefers \(c\) since \(2>1\).
Player 1 initially prefers \(a\) since \(3>2\).
All four comparisons are strict, establishing the full
backward contingent plan \((ae,cg)\) and its terminal outcome.

After the observed move \(b\), the strategies of player 1
consistent with the observation are \(be,bf\).
For any conjecture about player 2, \(be\) earns at most \(2\),
while either initial-\(a\) strategy guarantees \(3\).
Thus no such conjecture makes \(be\) originally optimal.
The strategy \(bf\), in contrast, earns \(4\) against \(dh\),
so it is optimal for that point-mass conjecture.
Consequently the stipulated rational-explanation filter
retains exactly \(bf\) after \(b\).

If player 2 expects the continuation \(f\) of that strategy,
then choosing \(d\) followed by the uniquely optimal \(g\)
earns \(4\), whereas \(c\) earns \(2\).
This gives the plan \(dg\). A first player who correctly
anticipates this plan gets \(3\) from \(ae\) or \(af\),
\(1\) from \(be\), and \(0\) from \(bf\), so chooses
\(a\) initially. This calculation distinguishes the
off-path strategy difference from the common outcome.
It does not assume that both conditional belief-revision
rules hold simultaneously at every hypothetical history.

For the quantitative rationalization, combine \(cg,ch\)
into their total probability \(p_c\). Let the other two
probabilities be \(p_{dg},p_{dh}\). Then \(bf\)'s value is
\(2p_c+4p_{dh}\); \(be\)'s is \(1+p_c\leq2\); and the
two initial-\(a\) values equal \(3\).
Hence \(bf\) is optimal exactly when
\(2p_c+4p_{dh}\geq3\).
Since \(p_{dh}\leq1-p_c\), this implies \(p_c\leq1/2\),
so continuation at \(d\) has positive probability.
Dividing the same inequality by \(4(1-p_c)\) gives
\[
 \frac{p_{dh}}{1-p_c}
 \geq\frac{3-2p_c}{4(1-p_c)}
 =\frac34+\frac{p_c}{4(1-p_c)}
 \geq\frac34 .
\]
This proves the precise high-\(h\)-probability implication,
including the weak-inequality and tie cases.
\end{proof}
\end{annexitem}
